% Svaki okvir daje jednu PDF stranu; korisnik je odobrio nastavak s010b.
\documentclass[aspectratio=169,11pt]{beamer}
\newcommand{\ETFTemaPutanja}{../_zajednicko/etf-v1}
\makeatletter
\edef\input@path{{\ETFTemaPutanja/}}
\makeatother
\usepackage{stil,dijagrami}
% Lokalno nastavno isticanje; zamrznuta tema etf-v1 ostaje neizmenjena.
\definecolor{DEisticanje}{HTML}{A32638}
\colorlet{DEcrvena}{DEisticanje}
\setbeamercolor{alerted text}{fg=DEisticanje}
\tikzset{DEstrelica/.style={wire,draw=DEisticanje,-{Latex[length=2mm]}}}
\renewcommand{\Oznaka}[1]{\par{\fontsize{10}{12}\selectfont\color{Muted}#1\par}}
\setlength{\parskip}{2pt}
\title[Digitalna elektronika 1 · 2021/22]{Kola srednjeg stepena integracije}
\author{prof. dr Lazar Saranovac}
\institute{Katedra za elektroniku, Elektrotehnički fakultet}
\date{2021/22}
\begin{document}
\slajdid{03-s001}
\begin{frame}[label=03-s001]
\centering\Huge
Kola\\[.5em]
srednjeg stepena integracije
\end{frame}

\slajdid{03-s002}
\begin{frame}[label=03-s002]{Integrisana logička kola i digitalni sistemi}
\gusttekst
Najčešća podela integrisanih logičkih kola i digitalnih sistema je na osnovu količine tranzistora koji se nalaze u istom kućištu, čipu. Ta podela je na SSI, MSI, LSI, VLSI i ULSI integrisana kola.

SSI - Small Scale Integration do 100 komponenti odnosno 10 gejtova

MSI - Medium Scale Integratiom do 500 komponeti odnosno od 10 do 100 gejtova

LSI – Large Scale Integration do 300000 komponenti odnosno više od 100 do 10000 gejtova

VLSI – Very Large Scale Integration više od 300000 komponenti 10000 do 100000 gejtova

ULSI – Ultra Large Scale Integration više od 1500000 komponenti preko 100000 gejtova

Ova podela po vrednostima nije striktna i u literaturi se mogu naći i drugačiji brojevi, pri čemu današnja terminologija pod pojmom VLSI dizajna podrazumeva i VLSI i ULSI tehnologije.
\end{frame}

\slajdid{03-s003}
\begin{frame}[label=03-s003]
Standardna pakovanja SSI tehnologije po pravilu imaju 14, 16 ili 20 nožica, pinova. Preko pinova se ostvaruje povezivanje putem veza na PCB (printed circuit board) sa ostalim komponentama. Uobičajeno je pin 7 (14 pinsko pakovanje), 8 (16 pinsko pakovanje) ili 10 (20 pinsko pakovanje) bio rezervisan za priključak mase, a pinovi 14, 16 i 20 za priključak napajanja. Ostali pinovi su rezervisani za ulaze i izlaze logičkih kola ili složenijih digitalnih sistema. Izbor tipa pakovanja zavisi od tehnologije u kojoj će PCB biti rađen. Proizvođači nude iste komponente u različitim pakovanjima i sve informacije o različitim dostupnim pakovanjima se nalaze u proizvođačkom opisu komponente.
\end{frame}

\slajdid{03-s004}
\begin{frame}[label=03-s004]
\gusttekst
Ono što je bitno jeste da su se u jednom pakovanju \underline{po pravilu} \underline{nalazila kola iste vrste,} \underline{na primer samo dvoulazna NI} \underline{kola, ili samo invertori itd\ldots} Znači ako nam za sintezu kombinacione mreže treba dva dvoulazna NI kola, jedno troulazno NI, i jedan invertor morali smo kupiti tri čipa: sa dvoulaznim NI kolima, sa troulaznim NI kolima i sa invertorima.

\alert{Ono što je takođe jako bitno jeste da se nekorišćeni ulazi i ulazi u nekorišćene delove komponente moraju postaviti na neaktivne nivoe. U suprotnom može doći do neželjenih efekata po rad aktivnog dela mreže. Kao što se vidi nije svejedno kako ćemo od početnih logičkih funkcija uraditi sintezu kombinacione mreže.}

U prethodnom primerima nadam se da ste uočili da sa dvoulaznim NI kolom možemo napraviti kratkim spajanjem ulaznih priključaka invertor – ne treba nam čip sa invertorima. Isto tako sa dva dvoulazna NI kola možemo napraviti jedno troulazno, normalno po cenu većeg kašnjenja \(F=A\cdot B\cdot C=A\cdot(B\cdot C)\), Znači trebalo bi nam i manji broj čipova, a i kupovali bi čipove iste vrste. Opet, ako smemo da dozvolimo, kašnjenje će biti povećano pošto umesto jednog kašnjenja kroz troulazno logičko kola imamo kašnjenje kroz dva dvoulazna logička kola.
\end{frame}

\slajdid{03-s005}
\begin{frame}[label=03-s005]{Dekoderi/demultiplekseri}
\gusttekst
Pod pojmom dekodera podrazumevaju se kombinacione mreže sa n ulaza i m izlaza pri čemu je \(n<m\). Najčešći su potpuni binarni dekoderi kod kojih je \(m=2^n\). Njihova struktura je u slučaju \(n=2\)
\begin{columns}[c,onlytextwidth]
\begin{column}{.52\textwidth}\centering
\begingroup
\def\DEcs{0}\def\DEokvir{0}
% Parametri: DEcs = 0/1; DEokvir = 0/1; DEnaziv, DEtip i DEulaz.
% Font se ne smanjuje skaliranjem kompletnog crteža.
\begin{tikzpicture}[circuit logic US,x=1cm,y=1cm,DEoznaka,
 inv/.style={not gate US,draw,minimum width=9mm,minimum height=8mm,scale=.7},
 kapija/.style={and gate US,draw,minimum width=10mm,minimum height=10mm,scale=.7,rotate=-90}]
\foreach \n/\y in {1/2.55,0/1.65}{
 \node[inv] (i\n) at (0,\y) {};
 \node[inv] (j\n) at (1.15,\y) {};
 \draw[DEvod] (i\n.input)--++(-.4,0) node[left] {S\n};
 \draw[DEvod] (i\n.output)--(j\n.input);
 \coordinate (tap\n) at ($(i\n.output)!.5!(j\n.input)$);
 \draw[DEvod] (j\n.output)--(5.25,\y);
 \draw[DEvod] (tap\n)--++(0,-.35)--(5.25,\y-.35);
}
\ifnum\DEcs=1
 \draw[DEvod] (-.95,.95) node[left] {\DEulaz} --(5.25,.95);
\fi
\foreach \n/\x/\sy/\ay in {0/2.05/2.2/1.3,1/3.05/2.2/1.65,2/4.05/2.55/1.3,3/5.05/2.55/1.65}{
 \ifnum\DEcs=1
  \node[kapija,logic gate inputs=nnn] (g\n) at (\x,.1) {};
  \draw[DEvod] (g\n.input 3)--(g\n.input 3 |- 0,\sy);
  \draw[DEvod] (g\n.input 2)--(g\n.input 2 |- 0,\ay);
  \draw[DEvod] (g\n.input 1)--(g\n.input 1 |- 0,.95);
 \else
  \node[kapija,logic gate inputs=nn] (g\n) at (\x,.1) {};
  \draw[DEvod] (g\n.input 2)--(g\n.input 2 |- 0,\sy);
  \draw[DEvod] (g\n.input 1)--(g\n.input 1 |- 0,\ay);
 \fi
 \draw[DEvod] (g\n.output)--++(0,-.25) node[below] {Y\n};
}
\ifnum\DEokvir=1
 \draw[DEvod,dashed] (-.75,-.42) rectangle (5.58,2.95);
 \node[align=center] at (.1,.3) {\DEnaziv\\\DEtip};
\fi
\end{tikzpicture}

\endgroup

\end{column}
\begin{column}{.45\textwidth}
Ulazi S se nazivaju selekcionim ulazima a često i adresnim ulazima sa oznakom A. Ono što je njihova osnovna karakteristika je da na svojim izlazima daju potpune proizvode formirane od ulaznih signala.
\(Y0=\overline{S1}\,\overline{S0}\), \(Y1=\overline{S1}\,S0\), \(Y2=S1\,\overline{S0}\) i \(Y3=S1\,S0\).
Prikazani dekoder naziva se potpunim binarnim dekoderom 2 u 4 (2/4). Pošto raspolažemo potpunim proizvodima formiranim od ulaznih signala sada možemo lako realizovati funkcije korišćenjem samo dodatnih ILI logičkih kola.
\end{column}
\end{columns}
\end{frame}

\slajdid{03-s006}
\begin{frame}[label=03-s006]
Na primer ako treba da realizujemo funkciju \(F=\bar B\bar A+BA\)
\par\centering
% Originalni slajd sadrži samo obris i oznake, bez unutrašnje šeme i izlazne veze.
\begin{tikzpicture}[DEoznaka,x=1cm,y=1cm]
\draw[DEvod,dashed] (0,0) rectangle (6,3);
\node[anchor=west] at (0,2.6) {S1};
\node[anchor=west] at (0,1.9) {S0};
\node[anchor=east] at (-.15,2.5) {B};
\node[anchor=east] at (-.15,1.8) {A};
\node[align=center] at (.9,.6) {Dekoder\\2/4};
\foreach \n/\x in {0/2.5,1/3.5,2/4.5,3/5.5}
 \node[anchor=north] at (\x,0) {Y\n};
\node at (4.4,-1.1) {F};
\end{tikzpicture}

\end{frame}

\slajdid{03-s007}
\begin{frame}[label=03-s007]
\gusttekst
Simboli mreža srednjeg stepena integracije su pravougaoni \underline{sa svim označenim signalima.}

Podvučeno je pošto na ispitu često zaboravite da označite signale unutar komponente, a onda nije moguće restaurirati šta ste i kako hteli da realizujete zadatak.
\begingroup
\def\DEcs{0}\def\DEokvir{1}\def\DEnaziv{Dekoder}\def\DEtip{2/4}
\begin{columns}[c,onlytextwidth]
\begin{column}{.62\textwidth}\centering
% Parametri: DEcs = 0/1; DEokvir = 0/1; DEnaziv, DEtip i DEulaz.
% Font se ne smanjuje skaliranjem kompletnog crteža.
\begin{tikzpicture}[circuit logic US,x=1cm,y=1cm,DEoznaka,
 inv/.style={not gate US,draw,minimum width=9mm,minimum height=8mm,scale=.7},
 kapija/.style={and gate US,draw,minimum width=10mm,minimum height=10mm,scale=.7,rotate=-90}]
\foreach \n/\y in {1/2.55,0/1.65}{
 \node[inv] (i\n) at (0,\y) {};
 \node[inv] (j\n) at (1.15,\y) {};
 \draw[DEvod] (i\n.input)--++(-.4,0) node[left] {S\n};
 \draw[DEvod] (i\n.output)--(j\n.input);
 \coordinate (tap\n) at ($(i\n.output)!.5!(j\n.input)$);
 \draw[DEvod] (j\n.output)--(5.25,\y);
 \draw[DEvod] (tap\n)--++(0,-.35)--(5.25,\y-.35);
}
\ifnum\DEcs=1
 \draw[DEvod] (-.95,.95) node[left] {\DEulaz} --(5.25,.95);
\fi
\foreach \n/\x/\sy/\ay in {0/2.05/2.2/1.3,1/3.05/2.2/1.65,2/4.05/2.55/1.3,3/5.05/2.55/1.65}{
 \ifnum\DEcs=1
  \node[kapija,logic gate inputs=nnn] (g\n) at (\x,.1) {};
  \draw[DEvod] (g\n.input 3)--(g\n.input 3 |- 0,\sy);
  \draw[DEvod] (g\n.input 2)--(g\n.input 2 |- 0,\ay);
  \draw[DEvod] (g\n.input 1)--(g\n.input 1 |- 0,.95);
 \else
  \node[kapija,logic gate inputs=nn] (g\n) at (\x,.1) {};
  \draw[DEvod] (g\n.input 2)--(g\n.input 2 |- 0,\sy);
  \draw[DEvod] (g\n.input 1)--(g\n.input 1 |- 0,\ay);
 \fi
 \draw[DEvod] (g\n.output)--++(0,-.25) node[below] {Y\n};
}
\ifnum\DEokvir=1
 \draw[DEvod,dashed] (-.75,-.42) rectangle (5.58,2.95);
 \node[align=center] at (.1,.3) {\DEnaziv\\\DEtip};
\fi
\end{tikzpicture}

\end{column}
\begin{column}{.33\textwidth}\centering
% Parametri DEcs, DEnaziv, DEtip, DEulaz.
\begin{tikzpicture}[DEoznaka,x=1cm,y=1cm]
\draw[DEvod] (0,0) rectangle (1.65,2.7);
\node[align=center] at (.82,2.25) {\DEnaziv\\\DEtip};
\foreach \n/\y in {3/1.72,2/1.32,1/.92,0/.52}{
 \draw[DEvod] (1.65,\y)--(1.95,\y);
 \node[anchor=east,inner sep=2pt] at (1.65,\y) {Y\n};
}
\foreach \n/\x in {1/.48,0/1.18}{
 \draw[DEvod] (\x,0)--(\x,-.3);
 \node[anchor=south,inner sep=2pt] at (\x,0) {S\n};
}
\ifnum\DEcs=1
 \draw[DEvod] (-.3,.52)--(0,.52);
 \node[anchor=west,inner sep=2pt] at (0,.52) {\DEulaz};
\fi
\end{tikzpicture}

\end{column}
\end{columns}
\endgroup
Selekcioni signali su indeksirani po težini, dok su indeksi izlaza odgovarajući indeksima proizvoda koji se formira na tom izlazu. I to je standardno i podrazumeva se.
\end{frame}

\slajdid{03-s008}
\begin{frame}[label=03-s008]
\gusttekst
Da bi se na lak način omogućilo pravljenje mreža većih kapaciteta, sa većim brojem ulaza i izlaza uobičajeno komponente imaju dodatne selekcione signale koji „selektuju“ komponentu, uobičajeno nazvani CS (chip select).
\begingroup
\def\DEcs{1}\def\DEokvir{1}\def\DEnaziv{Dekoder}\def\DEtip{2/4}\def\DEulaz{CS}
\begin{columns}[c,onlytextwidth]
\begin{column}{.57\textwidth}\centering
% Parametri: DEcs = 0/1; DEokvir = 0/1; DEnaziv, DEtip i DEulaz.
% Font se ne smanjuje skaliranjem kompletnog crteža.
\begin{tikzpicture}[circuit logic US,x=1cm,y=1cm,DEoznaka,
 inv/.style={not gate US,draw,minimum width=9mm,minimum height=8mm,scale=.7},
 kapija/.style={and gate US,draw,minimum width=10mm,minimum height=10mm,scale=.7,rotate=-90}]
\foreach \n/\y in {1/2.55,0/1.65}{
 \node[inv] (i\n) at (0,\y) {};
 \node[inv] (j\n) at (1.15,\y) {};
 \draw[DEvod] (i\n.input)--++(-.4,0) node[left] {S\n};
 \draw[DEvod] (i\n.output)--(j\n.input);
 \coordinate (tap\n) at ($(i\n.output)!.5!(j\n.input)$);
 \draw[DEvod] (j\n.output)--(5.25,\y);
 \draw[DEvod] (tap\n)--++(0,-.35)--(5.25,\y-.35);
}
\ifnum\DEcs=1
 \draw[DEvod] (-.95,.95) node[left] {\DEulaz} --(5.25,.95);
\fi
\foreach \n/\x/\sy/\ay in {0/2.05/2.2/1.3,1/3.05/2.2/1.65,2/4.05/2.55/1.3,3/5.05/2.55/1.65}{
 \ifnum\DEcs=1
  \node[kapija,logic gate inputs=nnn] (g\n) at (\x,.1) {};
  \draw[DEvod] (g\n.input 3)--(g\n.input 3 |- 0,\sy);
  \draw[DEvod] (g\n.input 2)--(g\n.input 2 |- 0,\ay);
  \draw[DEvod] (g\n.input 1)--(g\n.input 1 |- 0,.95);
 \else
  \node[kapija,logic gate inputs=nn] (g\n) at (\x,.1) {};
  \draw[DEvod] (g\n.input 2)--(g\n.input 2 |- 0,\sy);
  \draw[DEvod] (g\n.input 1)--(g\n.input 1 |- 0,\ay);
 \fi
 \draw[DEvod] (g\n.output)--++(0,-.25) node[below] {Y\n};
}
\ifnum\DEokvir=1
 \draw[DEvod,dashed] (-.75,-.42) rectangle (5.58,2.95);
 \node[align=center] at (.1,.3) {\DEnaziv\\\DEtip};
\fi
\end{tikzpicture}

\end{column}
\begin{column}{.40\textwidth}\centering
% Parametri DEcs, DEnaziv, DEtip, DEulaz.
\begin{tikzpicture}[DEoznaka,x=1cm,y=1cm]
\draw[DEvod] (0,0) rectangle (1.65,2.7);
\node[align=center] at (.82,2.25) {\DEnaziv\\\DEtip};
\foreach \n/\y in {3/1.72,2/1.32,1/.92,0/.52}{
 \draw[DEvod] (1.65,\y)--(1.95,\y);
 \node[anchor=east,inner sep=2pt] at (1.65,\y) {Y\n};
}
\foreach \n/\x in {1/.48,0/1.18}{
 \draw[DEvod] (\x,0)--(\x,-.3);
 \node[anchor=south,inner sep=2pt] at (\x,0) {S\n};
}
\ifnum\DEcs=1
 \draw[DEvod] (-.3,.52)--(0,.52);
 \node[anchor=west,inner sep=2pt] at (0,.52) {\DEulaz};
\fi
\end{tikzpicture}

\end{column}
\end{columns}
\endgroup
U ovom slučaju za dekoder se kaže da su mu aktivne logičke jedinice na izlazu, a CS signal dovodi sve izlaze u neaktivno stanje, odnosno kao što se vidi kada CS nije aktivan odnosno CS=0, svi izlazi će biti na neaktivnom nivou odnosno logičkoj nuli. Česta situacija je i da komponenta ima više od jednog CS signala. U tom slučaju su uvek povezani I logičkom funkcijom, odnosno svi moraju biti aktivni da bi se selektovala komponenta, \(CS=CS1\cdot CS2\).
\end{frame}

\slajdid{03-s009}
\begin{frame}[t,label=03-s009]
\gusttekst
\begin{columns}[c,onlytextwidth]
\begin{column}{.47\textwidth}\centering
\begingroup
% Ujednačen mali simbol sa 8 pt oznakama; selekcioni priključci su levo.
% \DEdek{ime}{x}{y}; globalno DEbar/DEbubble menjaju izlaze.
\providecommand{\DEbar}{0}
\providecommand{\DEbubble}{0}
\newcommand{\DEdek}[3]{%
\begin{scope}[shift={(#2,#3)},font=\fontsize{8}{8.5}\selectfont]
\draw[DEvod] (0,0) rectangle (1.7,1.65);
\node[align=center,inner sep=1pt] at (.58,1.27) {Dekoder\\2/4};
\foreach \n/\y in {3/1.24,2/.96,1/.68,0/.40}{
 \ifnum\DEbar=1
  \node[anchor=east,inner sep=1pt] at (1.7,\y) {$\overline{\mathrm{Y\n}}$};
 \else
  \node[anchor=east,inner sep=1pt] at (1.7,\y) {Y\n};
 \fi
 \ifnum\DEbubble=1
  \draw[DEvod,fill=white] (1.75,\y) circle(.05);
  \draw[DEvod] (1.80,\y)--(1.95,\y);
 \else
  \draw[DEvod] (1.7,\y)--(1.95,\y);
 \fi
 \coordinate (#1-y\n) at (1.95,\y);
}
\foreach \lab/\yy in {S1/.96,S0/.68,CS/.40}{
 \node[anchor=west,inner sep=1pt] at (0,\yy) {\lab};
 \draw[DEvod] (-.2,\yy)--(0,\yy);
 \coordinate (#1-\lab) at (-.2,\yy);
}
\end{scope}}

\begin{tikzpicture}[x=1cm,y=1cm,font=\fontsize{8}{9}\selectfont]
\draw[DEvod] (-.12,-.45) rectangle (4.76,7.3);
\node[align=center] at (.75,6.85) {Dekoder\\4/16};
\DEdek{r}{0}{3}
\foreach \i/\y in {3/5.5,2/3.7,1/1.9,0/.1}{
 \DEdek{d\i}{3}{\y}
 \foreach \j in {3,2,1,0}{
  \pgfmathtruncatemacro{\izlaz}{4*\i+\j}
  \draw[DEvod] (d\i-y\j)--++(.15,0) node[right,inner sep=1pt] {Y\izlaz};
 }
 \draw[DEvod] (d\i-S1)--(2.45,0 |- d\i-S1);
 \draw[DEvod] (d\i-S0)--(2.68,0 |- d\i-S0);
}
\draw[DEvod] (r-y3)--(2.1,4.24)|-(d3-CS);
\draw[DEvod] (r-y2)--(2.22,3.96)|-(d2-CS);
\draw[DEvod] (r-y1)--(2.1,3.68)|-(d1-CS);
\draw[DEvod] (r-y0)--(2.22,3.4)|-(d0-CS);
\draw[DEvod] (2.45,6.46)--(2.45,-.16)--(-.55,-.16) node[left] {S1};
\draw[DEvod] (2.68,6.18)--(2.68,-.35)--(-.55,-.35) node[left] {S0};
\draw[DEvod] (r-S1)--(-.55,3.96) node[left] {S3};
\draw[DEvod] (r-S0)--(-.55,3.68) node[left] {S2};
\draw[DEvod] (r-CS)--(-.55,3.40) node[left] {CS};
\end{tikzpicture}
\endgroup

\end{column}
\begin{column}{.50\textwidth}
Piramidalna struktura.

U ovom slučaju nam je dovoljan jedan selekcioni signal na komponentama

Mogući su i drugačiji načini povezivanja internih selekcionih signala. Prikazani način dozvoljava da se na lak način identifikuju indeksi novih izlaza. Po pravilu u izlaznom nivou se koriste selekcioni signali „najniže težine, a onda u prethodnim nivoima sve veće težine. Uočiti da na ovaj način možemo realizovati i komponentu sa većim brojem izlaza jedino će se povećavati broj nivoa. Na primer komponenta \(6/2^6\) realizovana sa dekoderima 2/4 bi imala tri nivoa. Naziv piramidalna struktura je zbog oblika koju ova struktura ima kao i načina dekodovanja.
\end{column}
\end{columns}
\end{frame}

\slajdid{03-s010}
\begin{frame}[t,label=03-s010]
\gusttekst
\fontsize{9}{9.5}\selectfont
\setlength{\parskip}{1pt}
Kod matrične strukture, treba nam dva selekciona signala na komponentama u slučaju dvodimenzionog dekodovanja, odnosno tru u slučaju trodimenzionog dekodovanja itd... Za veći broj izlaza matrična struktura će imati prednost zbog manjeg broja upotrebljenih komponenti, ali kao što smo rekli potrebne su specifičnije komponente. U tom smislu da vidimo kako bi izgledala realizacija dekodera \(6/2^6\) realizovana sa dekoderima 2/4.
\begin{columns}[c,onlytextwidth]
\begin{column}{.60\textwidth}\centering
\begingroup
\newcommand{\MDek}[3]{%
 \begin{scope}[shift={(#2,#3)}]
 \draw[DEvod] (0,0) rectangle(1.6,2.4);
 \node[align=center,inner sep=1pt] at (.8,1.90) {Dekoder\\2/4};
 \foreach \n/\y in {3/1.6,2/1.3,1/1.0,0/.7}{
  \node[anchor=east,inner sep=1pt] at (1.6,\y) {Y\n};
  \draw[DEvod] (1.6,\y)--(1.8,\y);
  \coordinate (#1-y\n) at (1.8,\y);
 }
 \foreach \n/\y in {1/1.0,2/.6}{
  \node[anchor=west,inner sep=1pt] at (0,\y) {CS\n};
  \draw[DEvod] (-.2,\y)--(0,\y);
  \coordinate (#1-cs\n) at (-.2,\y);
 }
 \foreach \n/\x in {1/.45,0/1.15}{
  \node[anchor=south,inner sep=1pt] at (\x,0) {S\n};
  \draw[DEvod] (\x,0)--(\x,-.2);
  \coordinate (#1-s\n) at (\x,-.2);
 }
 \end{scope}}
\begin{tikzpicture}[x=.86cm,y=.78cm,font=\fontsize{8}{9}\selectfont]
\MDek{v}{0}{2.8}
\node at (.27,4.2) {i};
\foreach \n/\y in {3/4.4,2/4.1,1/3.8,0/3.5}{
 \draw[DEvod] (v-y\n)--(3.75,\y);
 \node[anchor=south,inner sep=1pt] at (2.13,\y) {Y\n V};
}
\draw[DEvod] (1.0,0) rectangle (4.1,1.45);
\node[align=center,inner sep=1pt] at (1.70,.75) {Dekoder\\2/4};
\node at (3.4,.7) {j};
\foreach \n/\x in {3/2.65,2/3.0,1/3.35,0/3.7}{
 \draw[DEvod] (\x,1.45)--(\x,4.55);
 \node[rotate=90,anchor=east,inner sep=1pt] at (\x,1.45) {Y\n};
 \node[rotate=90,anchor=west,inner sep=1pt] at (\x,1.85) {Y\n K};
}
\foreach \n/\x in {1/2.55,2/3.0}{
 \node[rotate=90,anchor=west,inner sep=1pt] at (\x,0) {CS\n};
 \draw[DEvod] (\x,0)--(\x,-.12);
}
\node[rotate=90,anchor=south,inner sep=1pt] at (4.1,.40) {S1};
\node[rotate=90,anchor=south,inner sep=1pt] at (4.1,1.02) {S0};
\draw[DEvod] (4.1,.40)--(4.35,.40) node[right,inner sep=1pt] {S3};
\draw[DEvod] (4.1,1.02)--(4.35,1.02) node[right,inner sep=1pt] {S2};
\draw[DEvod] (v-cs1)--(-.35,3.8)--(-.35,-.22)--(2.55,-.22)--(2.55,-.12);
\draw[DEvod] (v-cs2)--(-.60,3.4)--(-.60,-.48)--(3.0,-.48)--(3.0,-.12);
\draw[DEvod] (-.35,-.22)--(-.8,-.22) node[left,inner sep=1pt] {CS1};
\draw[DEvod] (-.60,-.48)--(-.8,-.48) node[left,inner sep=1pt] {CS2};
\draw[DEvod] (v-s1)--++(0,-.25) node[below] {S5};
\draw[DEvod] (v-s0)--++(0,-.25) node[below] {S4};
\MDek{f}{4.8}{3}
\foreach \n in {3,2,1,0}
 \draw[DEvod] (f-y\n)--++(.08,0) node[right,inner sep=1pt] {Y(16i+4j+\n)};
\draw[DEvod] (f-cs1)--++(-.15,0) node[left,inner sep=1pt] {YiV};
\draw[DEvod] (f-cs2)--++(-.15,0) node[left,inner sep=1pt] {YjK};
\draw[DEvod] (f-s1)--++(0,-.25) node[below] {S1};
\draw[DEvod] (f-s0)--++(0,-.25) node[below] {S0};
\draw[DEvod,fill=white] (3.0,4.1) circle(.075);
\draw[DEstrelica] (3.0,4.18) .. controls (3.6,4.8) and (4.0,4.9) .. (4.5,4.9);
\end{tikzpicture}
\endgroup

\end{column}
\begin{column}{.37\textwidth}
U izlaznom nivou se nalazi 16 dekodera 2/4 (na slici je prikazan smo jedan) dok se jedan dekoder koristi za selekciju vrste (i) a drugi za selekciju kolone (j). Potreban broj komponenti nam je 16+1+1=18. Način povezivanja internih selekcionih signala je identičan kao I kod piramidalne strukture. Uočiti da bi nam kod piramidalne sturkture trebalo 16+4+1=21 komponenta.
\end{column}
\end{columns}

Ova razlika postaje još značajnija kao su u pitanju komponente sa još većim brojem izlaza. Takođe uočiti da će matrična struktura imati manje kašnjenje, pošto u ovom slučaju uvek su u pitanju dva nivoa, dok kod piramidalne strukture će broj nivoa zavisiti od veličine komponente. Ali za matričnu strukturu su potrebne komponenet sa više selekcionih signala. Normalno moguće je praviti i kombinacije ove dve strukture.

\end{frame}

\slajdid{03-s010b}
\begin{frame}[label=03-s010b]{Matrična i piramidalna struktura: poređenje}
% element: 03-s010b:e03
Za dekoder \(6/64\) realizovan dekoderima 2/4:
\par\vspace{2mm}
\begin{columns}[T,onlytextwidth]
\begin{column}{.475\textwidth}
\textbf{Matrična struktura}\par\vspace{1mm}
16 izlaznih dekodera + dekoder vrste + dekoder kolone
\[\textcolor{DEisticanje}{16+1+1=18\text{ komponenti}}\]
\end{column}
\begin{column}{.475\textwidth}
\textbf{Piramidalna struktura}\par\vspace{1mm}
16 dekodera u izlaznom + 4 u srednjem + 1 u prvom nivou
\[16+4+1=21\text{ komponenta}\]
\end{column}
\end{columns}
\par\vspace{2mm}
% element: 03-s010b:e04
\Rezultat{\textbf{Matrična struktura ima dva nivoa, pa i manje kašnjenje.}
U piramidalnoj strukturi broj nivoa zavisi od veličine dekodera;
za ovaj primer potrebna su tri nivoa.}
\par\vspace{2mm}
Pri većem broju izlaza prednost matrične strukture u broju komponenti postaje značajnija,
ali zahteva \textbf{specifičnije komponente sa više CS ulaza}.
Moguće je kombinovati matričnu i piramidalnu strukturu.
\end{frame}

\slajdid{03-s011}
\begin{frame}[label=03-s011]
\gusttekst
Takođe ulazi i izlazi mogu biti i sa aktivnim nivoima logičke nule. Na primer da su izlazi sa aktivnim logičkim nulama
\begin{columns}[c,onlytextwidth]
\begin{column}{.55\textwidth}\centering
\begingroup
% Ujednačen mali simbol sa 8 pt oznakama; selekcioni priključci su levo.
% \DEdek{ime}{x}{y}; globalno DEbar/DEbubble menjaju izlaze.
\providecommand{\DEbar}{0}
\providecommand{\DEbubble}{0}
\newcommand{\DEdek}[3]{%
\begin{scope}[shift={(#2,#3)},font=\fontsize{8}{8.5}\selectfont]
\draw[DEvod] (0,0) rectangle (1.7,1.65);
\node[align=center,inner sep=1pt] at (.58,1.27) {Dekoder\\2/4};
\foreach \n/\y in {3/1.24,2/.96,1/.68,0/.40}{
 \ifnum\DEbar=1
  \node[anchor=east,inner sep=1pt] at (1.7,\y) {$\overline{\mathrm{Y\n}}$};
 \else
  \node[anchor=east,inner sep=1pt] at (1.7,\y) {Y\n};
 \fi
 \ifnum\DEbubble=1
  \draw[DEvod,fill=white] (1.75,\y) circle(.05);
  \draw[DEvod] (1.80,\y)--(1.95,\y);
 \else
  \draw[DEvod] (1.7,\y)--(1.95,\y);
 \fi
 \coordinate (#1-y\n) at (1.95,\y);
}
\foreach \lab/\yy in {S1/.96,S0/.68,CS/.40}{
 \node[anchor=west,inner sep=1pt] at (0,\yy) {\lab};
 \draw[DEvod] (-.2,\yy)--(0,\yy);
 \coordinate (#1-\lab) at (-.2,\yy);
}
\end{scope}}

\begin{tikzpicture}[x=1cm,y=1cm,font=\fontsize{9}{10}\selectfont]
\def\DEbar{1}\def\DEbubble{0}\DEdek{a}{0}{0}
\node at (2.3,.83) {$=$};
\def\DEbar{0}\def\DEbubble{1}\DEdek{b}{2.9}{0}
\node at (5.2,.83) {$\ne$};
\def\DEbar{1}\def\DEbubble{1}\DEdek{c}{5.8}{0}
\end{tikzpicture}
\endgroup

\end{column}
\begin{column}{.42\textwidth}
Obratite pažnju da su i prvi i drugi način označavanja isti, i mogu slobodno da se koriste, dok je treći pogrešan i takva komponenta nije isto što i prve dve. Suštinski treća komponenta bi opet bila dekoder sa aktivnim logičkim jedinicama na izlazu, postoji dvostruka negacija. Način prikazan na drugom simbolu se i najčešće koristi.
\end{column}
\end{columns}
Primer sa više selekcionih signala.
\par\centering
\begin{tikzpicture}[x=1cm,y=1cm,font=\fontsize{8}{9}\selectfont]
\foreach \i/\x in {0/0,1/3,2/6}{
 \begin{scope}[shift={(\x,0)}]
 \draw[DEvod] (0,0) rectangle(1.65,2.25);
 \node[align=center] at (.82,1.9) {Dekoder\\2/4};
 \foreach \n/\y in {3/1.48,2/1.16,1/.84,0/.52}{
  \node[anchor=east,inner sep=2pt] at (1.65,\y) {Y\n};
  \draw[DEvod] (1.65,\y)--(1.85,\y);
 }
 \foreach \n/\xx in {1/.48,0/1.18}{
  \node[anchor=south,inner sep=2pt] at (\xx,0) {S\n};
  \draw[DEvod] (\xx,0)--(\xx,-.2);
 }
 \node[anchor=west,inner sep=2pt] at (0,.47) {CS1};
 \draw[DEvod] (-.2,.47)--(0,.47);
 \ifnum\i=1
  \node[anchor=west,inner sep=2pt] at (0,.82) {CS2};
 \else
  \node[anchor=west,inner sep=2pt] at (0,.82) {$\overline{\mathrm{CS2}}$};
 \fi
 \ifnum\i=0
  \draw[DEvod] (-.2,.82)--(0,.82);
 \else
  \draw[DEvod,fill=white] (-.05,.82) circle(.05);
  \draw[DEvod] (-.2,.82)--(-.1,.82);
 \fi
 \end{scope}
}
\node[font=\fontsize{11}{12}\selectfont] at (2.4,1.1) {$=$};
\node[font=\fontsize{11}{12}\selectfont] at (5.4,1.1) {$\ne$};
\end{tikzpicture}
\par
Unutrašnji selekcioni signal se formira kao \(CS=CS1\cdot\overline{CS2}\).
\end{frame}

\slajdid{03-s012}
\begin{frame}[label=03-s012]
\gusttekst
Demultipleksor je komponenta čija je osobina da sa ulaza (jednog) prosleđuje logički nivo na selektovani, adresirani, preko selekcionih signala, izlaz. Njegova unutrašnja struktura je
\par\centering
% Izvor prikazuje obrise/oznake bez unutrašnje šeme; desni simbol nema obris.
\begin{tikzpicture}[DEoznaka,x=1cm,y=1cm]
\draw[DEvod,dashed] (0,0) rectangle (4.5,2.4);
\foreach \lab/\y in {S1/2.1,S0/1.5,I/1.0}
 \node[anchor=east] at (-.15,\y) {\lab};
\node[align=center] at (.65,.55) {DMUX\\1/4};
\foreach \n/\x in {0/2.0,1/2.7,2/3.4,3/4.1}
 \node[anchor=north] at (\x,0) {Y\n};
\node[align=center] at (7.1,2.0) {DMUX\\1/4};
\foreach \n/\y in {3/1.8,2/1.4,1/1.0,0/.6}
 \node at (7.85,\y) {Y\n};
\node at (6.6,.55) {I};
\node at (6.9,.05) {S1};
\node at (7.5,.05) {S0};
\end{tikzpicture}
\par
\raggedright
Kao što se vidi, njegova unutrašnja struktura je identična dekoderu sa selekcionim signalom komponente. Zato se u katalozima dekoderi i demultipleksori nalaze pod zajedničkim imenom dekoder/demultiplekser. Način pravljenja mreža većih kapaciteta je identičan kako i kod dekoderskih mreža.
\end{frame}

\slajdid{03-s013}
\begin{frame}[label=03-s013]{Multiplekseri}
\gusttekst
Multiplekseri su specifične koderske mreže (broj ulaza veći od broja izlaza) koji imaju obrnutu funkciju od demultipleksera, odnosno ova komponenata prenosi stanje logičkog signala sa selektovanog, adresiranog, ulaza na izlaz. Znači ima n selekcionih signala, \(2^n\) ulaza i jedan izlaz.
\begingroup
\def\DEjedinica{.85cm}\def\DEdekmod{0}\def\DEcs{0}\def\DEprimer{0}
\begin{columns}[c,onlytextwidth]
\begin{column}{.64\textwidth}\centering
% DEdekmod=0: eksplicitni invertori; 1: blok dekodera. DEcs selekcija bloka.
\providecommand{\DEjedinica}{1cm}
\providecommand{\DEoblik}{.7}
\providecommand{\DEoznsize}{9}
\begin{tikzpicture}[circuit logic US,x=\DEjedinica,y=\DEjedinica,font=\fontsize{\DEoznsize}{\DEoznsize}\selectfont,
 inv/.style={not gate US,draw,minimum width=9mm,minimum height=8mm,scale=\DEoblik},
 kapija/.style={and gate US,draw,minimum width=10mm,minimum height=10mm,scale=\DEoblik,rotate=-90}]
\ifnum\DEdekmod=0
 \foreach \n/\y in {1/3.2,0/2.3}{
  \node[inv] (i\n) at (0,\y) {};
  \node[inv] (j\n) at (1.15,\y) {};
  \draw[DEvod] (i\n.input)--++(-.4,0) node[left] {S\n};
  \draw[DEvod] (i\n.output)--(j\n.input);
  \coordinate (tap\n) at ($(i\n.output)!.5!(j\n.input)$);
  \draw[DEvod] (j\n.output)--(5.4,\y);
  \draw[DEvod] (tap\n)--++(0,-.35)--(5.4,\y-.35);
 }
\else
 \draw[DEvod] (-.6,1.5) rectangle (1.2,3.85);
 \node[align=center] at (.3,3.4) {Dekoder\\2/4};
 \foreach \n/\y in {3/3.0,2/2.66,1/2.33,0/2.0}{
  \node[anchor=east,inner sep=1pt] at (1.2,\y) {Y\n};
  \draw[DEvod] (1.2,\y)--(5.4,\y);
 }
 \foreach \n/\x/\y in {1/-.18/1.23,0/.55/.96}{
  \node[anchor=south,inner sep=1pt] at (\x,1.5) {S\n};
  \draw[DEvod] (\x,1.5)--(\x,\y)--(-1,\y) node[left] {S\n};
 }
 \ifnum\DEcs=1
  \node[anchor=west,inner sep=1pt] at (-.6,2.03) {CS};
  \draw[DEvod] (-.6,2.03)--(-1,2.03) node[left] {CS};
 \fi
\fi
\ifnum\DEdekmod=1\def\DEvrh{4.2}\def\DEkutija{3.98}\else\def\DEvrh{3.98}\def\DEkutija{3.65}\fi
\foreach \n/\x/\sy/\ay/\dy in {0/2.05/2.85/1.95/2.0,1/3.05/2.85/2.3/2.33,2/4.05/3.2/1.95/2.66,3/5.05/3.2/2.3/3.0}{
 \ifnum\DEdekmod=0
  \node[kapija,logic gate inputs=nnn] (g\n) at (\x,.6) {};
  \draw[DEvod] (g\n.input 3)--(g\n.input 3 |- 0,\sy);
  \draw[DEvod] (g\n.input 2)--(g\n.input 2 |- 0,\ay);
 \else
  \node[kapija,logic gate inputs=nn] (g\n) at (\x,.6) {};
  \draw[DEvod] (g\n.input 2)--(g\n.input 2 |- 0,\dy);
 \fi
 \draw[DEvod] (g\n.input 1)--(g\n.input 1 |- 0,\DEvrh) node[above] {I\n};
}
\node[or gate US,draw,logic gate inputs=nnnn,minimum width=10mm,minimum height=12mm,scale=\DEoblik,rotate=-90] (z) at (3.55,-.92) {};
\foreach \n/\pin/\y in {0/4/-.22,1/3/0,2/2/0,3/1/-.22}
 \draw[DEvod] (g\n.output)--(g\n.output |- 0,\y)--(z.input \pin |- 0,\y)--(z.input \pin);
\draw[DEvod] (z.output)--++(0,-.25) node[below] {Z};
\coordinate (DEdno) at ($(z.output)+(0,-.08)$);
\draw[DEvod,dashed] (-.75,0 |- DEdno) rectangle (5.6,\DEkutija);
\node[align=center] at (.15,.15) {MUX\\4/1};
\end{tikzpicture}

\end{column}
\begin{column}{.30\textwidth}\centering
% DEcs: dodatni CS; DEprimer: spoljašnji nivoi 1,0,0,1 i B/A/F.
\begin{tikzpicture}[DEoznaka,x=1cm,y=1cm]
\ifnum\DEcs=1\def\DEpom{.35}\else\def\DEpom{0}\fi
\draw[DEvod] (0,0) rectangle (1.65,{2.7+\DEpom});
\node[align=center] at (.82,{2.25+\DEpom}) {MUX\\4/1};
\foreach \n/\y/\v in {3/1.72/1,2/1.32/0,1/.92/0,0/.52/1}{
 \draw[DEvod] (-.3,{\y+\DEpom})--(0,{\y+\DEpom});
 \node[anchor=west,inner sep=2pt] at (0,{\y+\DEpom}) {I\n};
 \ifnum\DEprimer=1
  \node[anchor=east] at (-.3,{\y+\DEpom}) {\v};
 \fi
}
\draw[DEvod] (1.65,{1.12+\DEpom})--(1.95,{1.12+\DEpom});
\node[anchor=east,inner sep=2pt] at (1.65,{1.12+\DEpom}) {Z};
\ifnum\DEprimer=1
 \node[anchor=west] at (1.95,{1.12+\DEpom}) {F};
\fi
\foreach \n/\x/\lab in {1/.48/B,0/1.18/A}{
 \draw[DEvod] (\x,0)--(\x,-.3);
 \node[anchor=south,inner sep=2pt] at (\x,0) {S\n};
 \ifnum\DEprimer=1
  \node[anchor=north] at (\x,-.3) {\lab};
 \fi
}
\ifnum\DEcs=1
 \draw[DEvod] (-.3,.48)--(0,.48);
 \node[anchor=west,inner sep=1pt] at (0,.48) {CS};
\fi
\end{tikzpicture}

\end{column}
\end{columns}
\endgroup
\end{frame}

\slajdid{03-s014}
\begin{frame}[label=03-s014]
\fontsize{9}{10}\selectfont
\setlength{\parskip}{2pt}
U nekim situacijama multiplekseri su komponente koje omogućuju još jednostavniju realizaciju kombinacionih mreža od dekodera. Na primer ako treba da realizujemo funkciju \(F=\bar B\bar A+BA\). Pre nego što prikažemo relizaciju ono što treba uočiti da je multiplekser principski
\begingroup
\def\DEoznsize{8}\def\DEoblik{.6}\def\DEdekmod{1}\def\DEcs{0}\def\DEprimer{0}\def\DEjedinica{.75cm}
\begin{columns}[c,onlytextwidth]
\begin{column}{.44\textwidth}\centering
% DEdekmod=0: eksplicitni invertori; 1: blok dekodera. DEcs selekcija bloka.
\providecommand{\DEjedinica}{1cm}
\providecommand{\DEoblik}{.7}
\providecommand{\DEoznsize}{9}
\begin{tikzpicture}[circuit logic US,x=\DEjedinica,y=\DEjedinica,font=\fontsize{\DEoznsize}{\DEoznsize}\selectfont,
 inv/.style={not gate US,draw,minimum width=9mm,minimum height=8mm,scale=\DEoblik},
 kapija/.style={and gate US,draw,minimum width=10mm,minimum height=10mm,scale=\DEoblik,rotate=-90}]
\ifnum\DEdekmod=0
 \foreach \n/\y in {1/3.2,0/2.3}{
  \node[inv] (i\n) at (0,\y) {};
  \node[inv] (j\n) at (1.15,\y) {};
  \draw[DEvod] (i\n.input)--++(-.4,0) node[left] {S\n};
  \draw[DEvod] (i\n.output)--(j\n.input);
  \coordinate (tap\n) at ($(i\n.output)!.5!(j\n.input)$);
  \draw[DEvod] (j\n.output)--(5.4,\y);
  \draw[DEvod] (tap\n)--++(0,-.35)--(5.4,\y-.35);
 }
\else
 \draw[DEvod] (-.6,1.5) rectangle (1.2,3.85);
 \node[align=center] at (.3,3.4) {Dekoder\\2/4};
 \foreach \n/\y in {3/3.0,2/2.66,1/2.33,0/2.0}{
  \node[anchor=east,inner sep=1pt] at (1.2,\y) {Y\n};
  \draw[DEvod] (1.2,\y)--(5.4,\y);
 }
 \foreach \n/\x/\y in {1/-.18/1.23,0/.55/.96}{
  \node[anchor=south,inner sep=1pt] at (\x,1.5) {S\n};
  \draw[DEvod] (\x,1.5)--(\x,\y)--(-1,\y) node[left] {S\n};
 }
 \ifnum\DEcs=1
  \node[anchor=west,inner sep=1pt] at (-.6,2.03) {CS};
  \draw[DEvod] (-.6,2.03)--(-1,2.03) node[left] {CS};
 \fi
\fi
\ifnum\DEdekmod=1\def\DEvrh{4.2}\def\DEkutija{3.98}\else\def\DEvrh{3.98}\def\DEkutija{3.65}\fi
\foreach \n/\x/\sy/\ay/\dy in {0/2.05/2.85/1.95/2.0,1/3.05/2.85/2.3/2.33,2/4.05/3.2/1.95/2.66,3/5.05/3.2/2.3/3.0}{
 \ifnum\DEdekmod=0
  \node[kapija,logic gate inputs=nnn] (g\n) at (\x,.6) {};
  \draw[DEvod] (g\n.input 3)--(g\n.input 3 |- 0,\sy);
  \draw[DEvod] (g\n.input 2)--(g\n.input 2 |- 0,\ay);
 \else
  \node[kapija,logic gate inputs=nn] (g\n) at (\x,.6) {};
  \draw[DEvod] (g\n.input 2)--(g\n.input 2 |- 0,\dy);
 \fi
 \draw[DEvod] (g\n.input 1)--(g\n.input 1 |- 0,\DEvrh) node[above] {I\n};
}
\node[or gate US,draw,logic gate inputs=nnnn,minimum width=10mm,minimum height=12mm,scale=\DEoblik,rotate=-90] (z) at (3.55,-.92) {};
\foreach \n/\pin/\y in {0/4/-.22,1/3/0,2/2/0,3/1/-.22}
 \draw[DEvod] (g\n.output)--(g\n.output |- 0,\y)--(z.input \pin |- 0,\y)--(z.input \pin);
\draw[DEvod] (z.output)--++(0,-.25) node[below] {Z};
\coordinate (DEdno) at ($(z.output)+(0,-.08)$);
\draw[DEvod,dashed] (-.75,0 |- DEdno) rectangle (5.6,\DEkutija);
\node[align=center] at (.15,.15) {MUX\\4/1};
\end{tikzpicture}

\end{column}
\begin{column}{.21\textwidth}\centering
% DEcs: dodatni CS; DEprimer: spoljašnji nivoi 1,0,0,1 i B/A/F.
\begin{tikzpicture}[DEoznaka,x=1cm,y=1cm]
\ifnum\DEcs=1\def\DEpom{.35}\else\def\DEpom{0}\fi
\draw[DEvod] (0,0) rectangle (1.65,{2.7+\DEpom});
\node[align=center] at (.82,{2.25+\DEpom}) {MUX\\4/1};
\foreach \n/\y/\v in {3/1.72/1,2/1.32/0,1/.92/0,0/.52/1}{
 \draw[DEvod] (-.3,{\y+\DEpom})--(0,{\y+\DEpom});
 \node[anchor=west,inner sep=2pt] at (0,{\y+\DEpom}) {I\n};
 \ifnum\DEprimer=1
  \node[anchor=east] at (-.3,{\y+\DEpom}) {\v};
 \fi
}
\draw[DEvod] (1.65,{1.12+\DEpom})--(1.95,{1.12+\DEpom});
\node[anchor=east,inner sep=2pt] at (1.65,{1.12+\DEpom}) {Z};
\ifnum\DEprimer=1
 \node[anchor=west] at (1.95,{1.12+\DEpom}) {F};
\fi
\foreach \n/\x/\lab in {1/.48/B,0/1.18/A}{
 \draw[DEvod] (\x,0)--(\x,-.3);
 \node[anchor=south,inner sep=2pt] at (\x,0) {S\n};
 \ifnum\DEprimer=1
  \node[anchor=north] at (\x,-.3) {\lab};
 \fi
}
\ifnum\DEcs=1
 \draw[DEvod] (-.3,.48)--(0,.48);
 \node[anchor=west,inner sep=1pt] at (0,.48) {CS};
\fi
\end{tikzpicture}

\end{column}
\begin{column}{.06\textwidth}\centering
\begin{tikzpicture}
\draw[DEstrelica,DEplava] (0,0)--(.6,0);
\draw[DEstrelica,DEplava] (0,1.6)--(.6,.9);
\end{tikzpicture}
\end{column}
\begin{column}{.26\textwidth}\centering
\def\DEprimer{1}
% DEcs: dodatni CS; DEprimer: spoljašnji nivoi 1,0,0,1 i B/A/F.
\begin{tikzpicture}[DEoznaka,x=1cm,y=1cm]
\ifnum\DEcs=1\def\DEpom{.35}\else\def\DEpom{0}\fi
\draw[DEvod] (0,0) rectangle (1.65,{2.7+\DEpom});
\node[align=center] at (.82,{2.25+\DEpom}) {MUX\\4/1};
\foreach \n/\y/\v in {3/1.72/1,2/1.32/0,1/.92/0,0/.52/1}{
 \draw[DEvod] (-.3,{\y+\DEpom})--(0,{\y+\DEpom});
 \node[anchor=west,inner sep=2pt] at (0,{\y+\DEpom}) {I\n};
 \ifnum\DEprimer=1
  \node[anchor=east] at (-.3,{\y+\DEpom}) {\v};
 \fi
}
\draw[DEvod] (1.65,{1.12+\DEpom})--(1.95,{1.12+\DEpom});
\node[anchor=east,inner sep=2pt] at (1.65,{1.12+\DEpom}) {Z};
\ifnum\DEprimer=1
 \node[anchor=west] at (1.95,{1.12+\DEpom}) {F};
\fi
\foreach \n/\x/\lab in {1/.48/B,0/1.18/A}{
 \draw[DEvod] (\x,0)--(\x,-.3);
 \node[anchor=south,inner sep=2pt] at (\x,0) {S\n};
 \ifnum\DEprimer=1
  \node[anchor=north] at (\x,-.3) {\lab};
 \fi
}
\ifnum\DEcs=1
 \draw[DEvod] (-.3,.48)--(0,.48);
 \node[anchor=west,inner sep=1pt] at (0,.48) {CS};
\fi
\end{tikzpicture}

\end{column}
\end{columns}
\endgroup
odnosno da ponovo imamo selekciju potpunih proizvoda koje formiraju signali S1 i S0. U tom smislu jedino što treba da uradimo jeste da na ulaze multipleksera za odgovarajući proizvod dovedemo željene vrednosti funkcije
\end{frame}

\slajdid{03-s015}
\begin{frame}[label=03-s015]
\gusttekst
Kao i kod dekoderskih mreža komponenta može imati selekcioni signal komponente
\begingroup
\def\DEoznsize{8}\def\DEoblik{.6}\def\DEdekmod{1}\def\DEcs{1}\def\DEprimer{0}\def\DEjedinica{.75cm}
\begin{columns}[c,onlytextwidth]
\begin{column}{.39\textwidth}\centering
% DEdekmod=0: eksplicitni invertori; 1: blok dekodera. DEcs selekcija bloka.
\providecommand{\DEjedinica}{1cm}
\providecommand{\DEoblik}{.7}
\providecommand{\DEoznsize}{9}
\begin{tikzpicture}[circuit logic US,x=\DEjedinica,y=\DEjedinica,font=\fontsize{\DEoznsize}{\DEoznsize}\selectfont,
 inv/.style={not gate US,draw,minimum width=9mm,minimum height=8mm,scale=\DEoblik},
 kapija/.style={and gate US,draw,minimum width=10mm,minimum height=10mm,scale=\DEoblik,rotate=-90}]
\ifnum\DEdekmod=0
 \foreach \n/\y in {1/3.2,0/2.3}{
  \node[inv] (i\n) at (0,\y) {};
  \node[inv] (j\n) at (1.15,\y) {};
  \draw[DEvod] (i\n.input)--++(-.4,0) node[left] {S\n};
  \draw[DEvod] (i\n.output)--(j\n.input);
  \coordinate (tap\n) at ($(i\n.output)!.5!(j\n.input)$);
  \draw[DEvod] (j\n.output)--(5.4,\y);
  \draw[DEvod] (tap\n)--++(0,-.35)--(5.4,\y-.35);
 }
\else
 \draw[DEvod] (-.6,1.5) rectangle (1.2,3.85);
 \node[align=center] at (.3,3.4) {Dekoder\\2/4};
 \foreach \n/\y in {3/3.0,2/2.66,1/2.33,0/2.0}{
  \node[anchor=east,inner sep=1pt] at (1.2,\y) {Y\n};
  \draw[DEvod] (1.2,\y)--(5.4,\y);
 }
 \foreach \n/\x/\y in {1/-.18/1.23,0/.55/.96}{
  \node[anchor=south,inner sep=1pt] at (\x,1.5) {S\n};
  \draw[DEvod] (\x,1.5)--(\x,\y)--(-1,\y) node[left] {S\n};
 }
 \ifnum\DEcs=1
  \node[anchor=west,inner sep=1pt] at (-.6,2.03) {CS};
  \draw[DEvod] (-.6,2.03)--(-1,2.03) node[left] {CS};
 \fi
\fi
\ifnum\DEdekmod=1\def\DEvrh{4.2}\def\DEkutija{3.98}\else\def\DEvrh{3.98}\def\DEkutija{3.65}\fi
\foreach \n/\x/\sy/\ay/\dy in {0/2.05/2.85/1.95/2.0,1/3.05/2.85/2.3/2.33,2/4.05/3.2/1.95/2.66,3/5.05/3.2/2.3/3.0}{
 \ifnum\DEdekmod=0
  \node[kapija,logic gate inputs=nnn] (g\n) at (\x,.6) {};
  \draw[DEvod] (g\n.input 3)--(g\n.input 3 |- 0,\sy);
  \draw[DEvod] (g\n.input 2)--(g\n.input 2 |- 0,\ay);
 \else
  \node[kapija,logic gate inputs=nn] (g\n) at (\x,.6) {};
  \draw[DEvod] (g\n.input 2)--(g\n.input 2 |- 0,\dy);
 \fi
 \draw[DEvod] (g\n.input 1)--(g\n.input 1 |- 0,\DEvrh) node[above] {I\n};
}
\node[or gate US,draw,logic gate inputs=nnnn,minimum width=10mm,minimum height=12mm,scale=\DEoblik,rotate=-90] (z) at (3.55,-.92) {};
\foreach \n/\pin/\y in {0/4/-.22,1/3/0,2/2/0,3/1/-.22}
 \draw[DEvod] (g\n.output)--(g\n.output |- 0,\y)--(z.input \pin |- 0,\y)--(z.input \pin);
\draw[DEvod] (z.output)--++(0,-.25) node[below] {Z};
\coordinate (DEdno) at ($(z.output)+(0,-.08)$);
\draw[DEvod,dashed] (-.75,0 |- DEdno) rectangle (5.6,\DEkutija);
\node[align=center] at (.15,.15) {MUX\\4/1};
\end{tikzpicture}

\end{column}
\begin{column}{.16\textwidth}\centering
% DEcs: dodatni CS; DEprimer: spoljašnji nivoi 1,0,0,1 i B/A/F.
\begin{tikzpicture}[DEoznaka,x=1cm,y=1cm]
\ifnum\DEcs=1\def\DEpom{.35}\else\def\DEpom{0}\fi
\draw[DEvod] (0,0) rectangle (1.65,{2.7+\DEpom});
\node[align=center] at (.82,{2.25+\DEpom}) {MUX\\4/1};
\foreach \n/\y/\v in {3/1.72/1,2/1.32/0,1/.92/0,0/.52/1}{
 \draw[DEvod] (-.3,{\y+\DEpom})--(0,{\y+\DEpom});
 \node[anchor=west,inner sep=2pt] at (0,{\y+\DEpom}) {I\n};
 \ifnum\DEprimer=1
  \node[anchor=east] at (-.3,{\y+\DEpom}) {\v};
 \fi
}
\draw[DEvod] (1.65,{1.12+\DEpom})--(1.95,{1.12+\DEpom});
\node[anchor=east,inner sep=2pt] at (1.65,{1.12+\DEpom}) {Z};
\ifnum\DEprimer=1
 \node[anchor=west] at (1.95,{1.12+\DEpom}) {F};
\fi
\foreach \n/\x/\lab in {1/.48/B,0/1.18/A}{
 \draw[DEvod] (\x,0)--(\x,-.3);
 \node[anchor=south,inner sep=2pt] at (\x,0) {S\n};
 \ifnum\DEprimer=1
  \node[anchor=north] at (\x,-.3) {\lab};
 \fi
}
\ifnum\DEcs=1
 \draw[DEvod] (-.3,.48)--(0,.48);
 \node[anchor=west,inner sep=1pt] at (0,.48) {CS};
\fi
\end{tikzpicture}
\par
\begin{tikzpicture}\draw[DEstrelica,DEplava] (0,0)--(1.4,0);\end{tikzpicture}
\end{column}
\begin{column}{.42\textwidth}\centering
\begingroup
\newcommand{\MMux}[3]{%
 \begin{scope}[shift={(#2,#3)}]
 \draw[DEvod] (0,0) rectangle (1.65,1.6);
 \node[align=center,inner sep=1pt] at (.9,1.18) {MUX\\4/1};
 \foreach \n/\y in {3/1.04,2/.78,1/.52,0/.26}{
  \node[anchor=west,inner sep=1pt] at (0,\y) {I\n};
  \draw[DEvod] (-.2,\y)--(0,\y);
  \coordinate (#1-i\n) at (-.2,\y);
 }
 \node[anchor=east,inner sep=1pt] at (1.65,.72) {Z};
 \draw[DEvod] (1.65,.72)--(1.9,.72);
 \coordinate (#1-z) at (1.9,.72);
 \foreach \n/\x in {1/.65,0/1.25}{
  \node[anchor=south,inner sep=1pt] at (\x,0) {S\n};
  \draw[DEvod] (\x,0)--(\x,-.1);
  \coordinate (#1-s\n) at (\x,-.1);
 }
 \end{scope}}
\begin{tikzpicture}[x=.85cm,y=.85cm,font=\fontsize{8}{8.5}\selectfont]
\foreach \k/\y in {3/5.85,2/3.90,1/1.95,0/0}{
 \MMux{m\k}{0}{\y}
 \foreach \n in {3,2,1,0}{
  \pgfmathtruncatemacro{\ii}{4*\k+\n}
  \draw[DEvod] (m\k-i\n)--++(-.15,0) node[left,inner sep=1pt] {I\ii};
 }
 \draw[DEvod] (m\k-s1)--(.65,\y-.13)--(2.05,\y-.13);
 \draw[DEvod] (m\k-s0)--(1.25,\y-.26)--(2.26,\y-.26);
}
\MMux{f}{3.35}{2.85}
\foreach \k/\x in {3/2.45,2/2.65,1/2.65,0/2.45}
 \draw[DEvod] (m\k-z)--(\x,0 |- m\k-z)|-(f-i\k);
\draw[DEvod] (2.05,5.72)--(2.05,-.5) node[below left,inner sep=1pt] {S1};
\draw[DEvod] (2.26,5.59)--(2.26,-.5) node[below right,inner sep=1pt] {S0};
\draw[DEvod] (f-z)--++(.2,0) node[right] {Z};
\draw[DEvod] (f-s1)--++(0,-.25) node[below] {S3};
\draw[DEvod] (f-s0)--++(0,-.25) node[below] {S2};
\end{tikzpicture}
\endgroup

\end{column}
\end{columns}
\endgroup
\end{frame}

\slajdid{03-s016}
\begin{frame}[label=03-s016]{Koder prioriteta}
\gusttekst
Osim multipleksera koji prestavlja specijalnu vrstu kodera u sintezi složenih digitalnih sistema (pogotovo u procesorskim) se najčešće koristi koder prioriteta. Osnovna njegova uloga jeste da koduje indeks ulaza na kojem se nalazi aktivan signal i najveće je vrednosti. Da bi bilo jasnije da vidimo funkcionalnu tabelu kodera prioriteta koji ima 4 ulaza. Da bi se kodovali indeksi potrebne su nam dve izlazne promenjive.
\begin{columns}[c,onlytextwidth]
\begin{column}{.41\textwidth}
\begingroup
\setlength{\tabcolsep}{3pt}\renewcommand{\arraystretch}{1.12}
\par\noindent
\begin{tabularx}{\linewidth}{*{6}{>{\centering\arraybackslash}X}}
\toprule I3&I2&I1&I0&A1&A0\\\midrule
1&X&X&X&1&1\\\cmidrule(lr){1-6}
0&1&X&X&1&0\\\cmidrule(lr){1-6}
0&0&1&X&0&1\\\cmidrule(lr){1-6}
0&0&0&1&0&0\\\cmidrule(lr){1-6}
0&0&0&0&0&0\\\bottomrule
\end{tabularx}\par
\endgroup

\end{column}
\begin{column}{.55\textwidth}
Neodređenost se pojavljuje u dve poslednje vrste. Pretposlednja vrsta je situacija da je na ulazu sa najnižim indeksom aktivan signal, a poslednja da uopšte nema aktivnih signala. Iz tog razloga nam treba još jedan dodatni izlazni signal GS (group select) koji će pokazivati da li uopšte ima nešto da se koduje, odnosno da li postoji bilo koji aktivan ulazni signal
\end{column}
\end{columns}
\end{frame}

\slajdid{03-s017}
\begin{frame}[label=03-s017]
\gusttekst
Proširena funkcionalna tabela u tom slučaju je
\begin{columns}[c,onlytextwidth]
\begin{column}{.35\textwidth}
\begingroup
\setlength{\tabcolsep}{3pt}\renewcommand{\arraystretch}{1.3}
\par\noindent
\begin{tabularx}{\linewidth}{*{7}{>{\centering\arraybackslash}X}}
\toprule I3&I2&I1&I0&A1&A0&GS\\\midrule
1&X&X&X&1&1&1\\\cmidrule(lr){1-7}
0&1&X&X&1&0&1\\\cmidrule(lr){1-7}
0&0&1&X&0&1&1\\\cmidrule(lr){1-7}
0&0&0&1&0&0&1\\\cmidrule(lr){1-7}
0&0&0&0&0&0&0\\\bottomrule
\end{tabularx}\par
\endgroup


Mada je moguće izvesti minimizaciju, u ovom slučaju je dosta jednostavnije direktno nacrtati realizaciju, ne trudeći se da bude baš minimalna. Direktno crtanje je moguće pošto su funkcije dosta „pravolinijske“.
\end{column}
\begin{column}{.62\textwidth}\centering
\begingroup
\def\DEenable{0}\def\DEjedinica{.9cm}\def\DEgejt{.65}
% DEenable: dodatni EI/EO; DEjedinica i DEgejt određuju geometriju, font ostaje 8 pt.
\begin{tikzpicture}[circuit logic US,x=\DEjedinica,y=\DEjedinica,font=\fontsize{8}{9}\selectfont,
 inv/.style={not gate US,draw,minimum width=8mm,minimum height=7mm,scale=\DEgejt},
 kapija/.style={and gate US,draw,minimum width=10mm,minimum height=10mm,scale=\DEgejt},
 zbir/.style={or gate US,draw,minimum width=10mm,minimum height=12mm,scale=\DEgejt,rotate=-90}]
\node[kapija,logic gate inputs=nn,anchor=input 2] (p2) at (2.7,3.2) {};
\node[kapija,logic gate inputs=nnn,anchor=input 3] (p1) at (2.7,2) {};
\node[kapija,logic gate inputs=nnnn,anchor=input 4] (p0) at (2.7,.8) {};
\draw[DEvod] (-.4,4.4) node[left] {I3} -- (3.4,4.4);
\coordinate (p3) at (3.4,4.4);
\foreach \n/\y/\pin in {2/3.2/2,1/2/3,0/.8/4}
 \draw[DEvod] (-.4,\y) node[left] {I\n} --(p\n.input \pin);
\foreach \n/\x/\y/\src in {3/.6/3.9/4.4,2/1.1/2.7/3.2,1/1.6/1.5/2}{
 \node[inv,rotate=-90] (i\n) at (\x,\y) {};
 \draw[DEvod] (\x,\src)--(i\n.input);
}
\foreach \n in {2,1,0}
 \draw[DEvod] (i3.output)|-(p\n.input 1);
\foreach \n in {1,0}
 \draw[DEvod] (i2.output)|-(p\n.input 2);
\draw[DEvod] (i1.output)|-(p0.input 3);
\foreach \n in {2,1,0} \coordinate (p\n) at (p\n.output);
\ifnum\DEenable=1
 \foreach \n in {3,2,1,0}{
  \node[kapija,logic gate inputs=nn,anchor=input 2] (e\n) at (4.5,0 |- p\n) {};
  \draw[DEvod] (p\n)--(e\n.input 2);
  \draw[DEvod] (e\n.input 1)--(4.03,0 |- e\n.input 1);
  \coordinate (q\n) at (e\n.output);
 }
 \draw[DEvod] (-.4,5.05) node[left] {EI} --(4.03,5.05)--(4.03,0 |- e0.input 1);
 \def\DXone{6.4}\def\DXzero{7.6}\def\DXgs{8.9}\def\DXend{9.45}
\else
 \foreach \n in {3,2,1,0} \coordinate (q\n) at (p\n);
 \def\DXone{4.5}\def\DXzero{5.7}\def\DXgs{7.1}\def\DXend{7.7}
 \node at (1.45,5.05) {Prioritetna mreža};
 \node at (5.9,5.05) {Koderska mreža};
\fi
\foreach \n in {3,2,1,0} \draw[DEvod] (q\n)--(\DXend,0 |- q\n);
\node[zbir,logic gate inputs=nn] (a1) at (\DXone,-.25) {};
\node[zbir,logic gate inputs=nn] (a0) at (\DXzero,-.25) {};
\node[zbir,logic gate inputs=nnnn] (gs) at (\DXgs,-.25) {};
\foreach \pin/\n in {2/3,1/2}
 \draw[DEvod] (a1.input \pin)--(a1.input \pin |- q\n);
\foreach \pin/\n in {2/3,1/1}
 \draw[DEvod] (a0.input \pin)--(a0.input \pin |- q\n);
\foreach \pin/\n in {4/3,3/2,2/1,1/0}
 \draw[DEvod] (gs.input \pin)--(gs.input \pin |- q\n);
\foreach \name/\lab in {a1/A1,a0/A0,gs/GS}
 \draw[DEvod] (\name.output)--(\name.output |- 0,-1.3) node[below] {\lab};
\ifnum\DEenable=1
 \node[inv,rotate=90] (ngs) at (10,-.35) {};
 \draw[DEvod] (gs.output)--(gs.output |- 0,-1.1)--(10,-1.1)--(ngs.input);
 \node[kapija,rotate=-90,logic gate inputs=nn] (eo) at (11.15,-.35) {};
 \draw[DEvod] (ngs.output)--(ngs.output |- 0,.35)--(eo.input 2 |- 0,.35)--(eo.input 2);
 \draw[DEvod] (4.03,5.05)--(4.03,5.5)--(11.6,5.5)--(11.6,.35)--(eo.input 1 |- 0,.35)--(eo.input 1);
 \draw[DEvod] (eo.output)--(eo.output |- 0,-1.3) node[below] {EO};
\fi
\end{tikzpicture}

\endgroup
\end{column}
\end{columns}
\end{frame}

\slajdid{03-s018}
\begin{frame}[label=03-s018]
\gusttekst
Prioritetna mreža će obezbediti da se u unutrašnjosti komponente pojavi samo jedna aktivna logička jedinica sa onog ulaza koji ima aktivnu jedinicu i najvišeg je prioriteta odnosno ima najveći indeks. U tom slučaju je lako napraviti Koderski deo mreže. Može i minimalnije da se nacrta
\par\centering
\begin{tikzpicture}[circuit logic US,DEoznaka,
 i/.style={not gate US,draw,minimum width=9mm,minimum height=8mm,scale=.7},
 g/.style={draw,minimum width=12mm,minimum height=10mm,scale=.7}]
\node[or gate US,g,logic gate inputs=nn] (a1) at (4.2,3.1) {};
\node[or gate US,g,logic gate inputs=nn] (a0) at (4.2,1.6) {};
\node[or gate US,g,logic gate inputs=nnnn] (gs) at (4.2,0) {};
\node[and gate US,g,logic gate inputs=nn] (p) at (2.55,1.35) {};
\node[i,anchor=output] (n) at (1.75,0 |- p.input 1) {};
\draw[DEvod] (a1.input 1)--(-.5,0 |- a1.input 1) node[left] {I3};
\draw[DEvod] (a1.input 2)--(-.5,0 |- a1.input 2) node[left] {I2};
\draw[DEvod] (p.input 2)--(-.5,0 |- p.input 2) node[left] {I1};
\draw[DEvod] (gs.input 4)--(-.5,0 |- gs.input 4) node[left] {I0};
\draw[DEvod] (.10,0 |- a1.input 2)|-(n.input);
\draw[DEvod] (n.output)--(p.input 1);
\draw[DEvod] (p.output)--++(.25,0)|-(a0.input 2);
\draw[DEvod] (3.45,0 |- a1.input 1)|-(a0.input 1);
\draw[DEvod] (.10,0 |- a1.input 2)|-(gs.input 1);
\draw[DEvod] (-.08,0 |- a1.input 1)|-(gs.input 2);
\draw[DEvod] (-.27,0 |- p.input 2)|-(gs.input 3);
\foreach \n/\lab in {a1/A1,a0/A0,gs/GS}
 \draw[DEvod] (\n.output)--++(.3,0) node[right] {\lab};
\end{tikzpicture}
\par
\raggedright
međutim prva realizacija je lakša za proširivanje kada je potrebno realizovati veće mreže.
\end{frame}

\slajdid{03-s019}
\begin{frame}[label=03-s019]
\gusttekst
Prilikom pravljenja mreža većih kapaciteta, možemo se poslužiti idejama koje smo imali kod prethodnih kola. Na primer da je potrebno napraviti koder prioriteta sa 16 ulaza
\begin{columns}[c,onlytextwidth]
\begin{column}{.43\textwidth}\centering
\begingroup
\def\DEenable{0}
% \KPSimbol{ime}{x}{y}; DEenable bira EI/EO varijantu; font 8 pt u pozivaocu.
\newcommand{\KPSimbol}[3]{%
\begin{scope}[shift={(#2,#3)}]
\draw[DEvod] (0,0) rectangle (2.1,1.65);
\ifnum\DEenable=1
 \def\DEgsheight{1.05}
 \node at (1.0,1.05) {KP};
 \node[anchor=north,inner sep=1pt] at (1.0,1.65) {EI};
 \draw[DEvod] (1.0,1.65)--(1.0,1.8);
 \coordinate (#1-ei) at (1.0,1.8);
 \node[anchor=south,inner sep=1pt] at (1.0,0) {EO};
 \draw[DEvod] (1.0,0)--(1.0,-.15);
 \coordinate (#1-eo) at (1.0,-.15);
 \foreach \n/\y in {1/.62,0/.25}{
  \node[anchor=east,inner sep=1pt] at (2.1,\y) {A\n};
  \draw[DEvod] (2.1,\y)--(2.3,\y);
  \coordinate (#1-a\n) at (2.3,\y);
 }
\else
 \def\DEgsheight{.95}
 \node at (1.0,1.43) {KP};
 \foreach \n/\x in {1/.7,0/1.45}{
  \node[anchor=south,inner sep=1pt] at (\x,0) {A\n};
  \draw[DEvod] (\x,0)--(\x,-.15);
  \coordinate (#1-a\n) at (\x,-.15);
 }
\fi
\foreach \n/\y in {3/1.2,2/.89,1/.58,0/.27}{
 \node[anchor=west,inner sep=1pt] at (0,\y) {I\n};
 \draw[DEvod] (-.2,\y)--(0,\y);
 \coordinate (#1-i\n) at (-.2,\y);
}
\node[anchor=east,inner sep=1pt] at (2.1,\DEgsheight) {GS};
\draw[DEvod] (2.1,\DEgsheight)--(2.3,\DEgsheight);
\coordinate (#1-gs) at (2.3,\DEgsheight);
\end{scope}}

\begin{tikzpicture}[x=.8cm,y=.8cm,font=\fontsize{8}{9}\selectfont]
\foreach \k/\y in {3/6,2/4,1/2,0/0}{
 \KPSimbol{p\k}{0}{\y}
 \foreach \n in {3,2,1,0}{
  \pgfmathtruncatemacro{\ii}{4*\k+\n}
  \draw[DEvod] (p\k-i\n)--++(-.15,0) node[left,inner sep=1pt] {I\ii};
 }
}
\KPSimbol{f}{3.5}{3}
\foreach \k/\x in {3/3.0,2/2.7,1/2.7,0/3.0}
 \draw[DEvod] (p\k-gs)--(\x,0 |- p\k-gs)|-(f-i\k);
\draw[DEvod] (f-gs)--++(.2,0) node[right] {GS};
\draw[DEvod] (f-a1)--++(0,-.2) node[below] {A3};
\draw[DEvod] (f-a0)--++(0,-.2) node[below] {A2};
\end{tikzpicture}
\endgroup

\end{column}
\begin{column}{.54\textwidth}
Sa slike je vidljiva ideja da se iskoristi osobina kodera prioriteta i u drugom nivou. Znači ako ima na bilo kojem koderu prioriteta u prvom nivou aktivnih logičkih jedinica, biće aktivan i njegov izlazni signal GS, pa će koder prioriteta u drugom nivou dati odgovarajući indeks kodera prioriteta prvog nivoa (određen po prioritetu GS signala), što su u stvari signali A3 i A2 za „veliki“ koder. Ostaje problem kako odrediti signale A1 i A0. \alert{Ovde su u pitanju izlazi pa njihovo direktno spajanje sigurno ne dolazi u obzir.} A i logički njihovo direktno spajanje ne bi donelo validnu informaciju. U tom smislu bi morali na osnovne kodere prioriteta dodati još neke ulaze i izlaze, kako bi omogućili laku identifikaciju koji signali A1 i A0 su validni, a i da omogućimo njihovo lako logičko spajanje.
\end{column}
\end{columns}
\end{frame}

% Original: PDF strana 10, donji slajd; prvi deo.
\slajdid{03-s020}
\begin{frame}[t,label=03-s020]{Ulaz EI i izlaz EO kodera prioriteta}
\setlength{\parskip}{0pt}
\fontsize{11}{13}\selectfont
% element: 03-s020:e01
Ulaz \textcolor{DEisticanje}{\textbf{EI}} (\emph{enable input}) daje validne izlaze \(A_1,A_0\)
kada je aktivan, a postavlja ih na \textbf{logičke nule} kada je neaktivan.
I \(GS\) možemo usloviti: njegova nula znači da nema aktivnih ulaza ili da je
\textbf{kodovanje zabranjeno} zato što koder višeg prioriteta ima šta da koduje.
\par\vspace{.5mm}
\begin{columns}[c,onlytextwidth]
\begin{column}{.67\textwidth}\centering
% element: 03-s020:e02
\begin{tikzpicture}[circuit logic US,x=1mm,y=.82mm,font=\fontsize{9}{11}\selectfont,
 kapija/.style={and gate US,draw,line width=.8pt,minimum width=10mm,minimum height=10mm,scale=.45},
 zbir/.style={or gate US,draw,line width=.8pt,minimum width=10mm,minimum height=12mm,scale=.45,rotate=-90}]
\node[kapija,logic gate inputs=nn,anchor=input 2] (p2) at (23,21) {};
\node[kapija,logic gate inputs=nnn,anchor=input 3] (p1) at (23,13) {};
\node[kapija,logic gate inputs=nnnn,anchor=input 4] (p0) at (23,5) {};
\draw[DEvod] (-3,29) node[left,inner sep=1pt] {I3}--(29,29);
\coordinate (p3) at (29,29);
\foreach \n/\yy/\pin in {2/21/2,1/13/3,0/5/4}
 \draw[DEvod] (-3,\yy) node[left,inner sep=1pt] {I\n}--(p\n.input \pin);
\foreach \n/\xx/\top/\source/\bus in {3/9/27.9/29/15.5,2/5/19.9/21/12.5,1/1/11.9/13/9.5}{
 \draw[DEvod] (\xx,\source)--(\xx,\top);
 \node[junction] at (\xx,\source) {};
 \draw[DEvod] (\xx-2,\top)--(\xx+2,\top)--(\xx,\top-3.2)--cycle;
 \draw[DEvod] (\xx,\top-3.55) circle[radius=.287mm];
 \coordinate (b\n) at (\bus,\top-3.9);
 \draw[DEvod] (\xx,\top-3.9)--(b\n);
}
\foreach \n in {2,1,0} \draw[DEvod] (b3)|-(p\n.input 1);
\foreach \n in {1,0} \draw[DEvod] (b2)|-(p\n.input 2);
\draw[DEvod] (b1)|-(p0.input 3);
\node[junction] at (b3 |- p2.input 1) {};
\node[junction] at (b3 |- p1.input 1) {};
\node[junction] at (b2 |- p1.input 2) {};
\foreach \n in {2,1,0} \coordinate (p\n) at (p\n.output);
\foreach \n in {3,2,1,0}{
 \node[kapija,logic gate inputs=nn,anchor=input 2] (e\n) at (38,0 |- p\n) {};
 \draw[DEvod] (p\n)--(e\n.input 2);
 \draw[DEvod] (e\n.input 1)--(35.5,0 |- e\n.input 1);
 \node[junction] at (35.5,0 |- e\n.input 1) {};
 \coordinate (q\n) at (e\n.output);
}
\draw[DEvod] (-3,36) node[left,inner sep=1pt] {EI}--(35.5,36)--(35.5,0 |- e0.input 1);
\node[junction] at (35.5,36) {};
\foreach \n in {3,2,1,0} \draw[DEvod] (q\n)--(74.5,0 |- q\n);
\node[zbir,logic gate inputs=nn] (a1) at (52,-1.5) {};
\node[zbir,logic gate inputs=nn] (a0) at (61,-1.5) {};
\node[zbir,logic gate inputs=nnnn] (gs) at (71,-1.5) {};
\foreach \pin/\n in {2/3,1/2}{
 \draw[DEvod] (a1.input \pin)--(a1.input \pin |- q\n);
 \node[junction] at (a1.input \pin |- q\n) {};
}
\foreach \pin/\n in {2/3,1/1}{
 \draw[DEvod] (a0.input \pin)--(a0.input \pin |- q\n);
 \node[junction] at (a0.input \pin |- q\n) {};
}
\foreach \pin/\n in {4/3,3/2,2/1,1/0}{
 \draw[DEvod] (gs.input \pin)--(gs.input \pin |- q\n);
 \node[junction] at (gs.input \pin |- q\n) {};
}
\foreach \name/\lab in {a1/A1,a0/A0}
 \draw[DEvod] (\name.output)--(\name.output |- 0,-12) node[below,inner sep=1pt] {\lab};
\coordinate (gs-spoj) at (gs.output |- 0,-8.5);
\draw[DEvod] (gs.output)--(gs-spoj)--(gs.output |- 0,-12) node[below,inner sep=1pt] {GS};
\node[junction] at (gs-spoj) {};
\draw[DEvod] (77,-3.2)--(81,-3.2)--(79,0)--cycle;
\draw[DEvod] (79,.35) circle[radius=.287mm];
\draw[DEvod] (gs-spoj)--(79,-8.5)--(79,-3.2);
\node[kapija,rotate=-90,logic gate inputs=nn] (eo) at (85.5,-1.5) {};
\draw[DEvod] (79,.7)--(79,3)--(eo.input 2 |- 0,3)--(eo.input 2);
\draw[DEvod] (35.5,36)--(35.5,38)--(90,38)--(90,3)--(eo.input 1 |- 0,3)--(eo.input 1);
\draw[DEvod] (eo.output)--(eo.output |- 0,-12) node[below,inner sep=1pt] {EO};
\end{tikzpicture}

\end{column}
\begin{column}{.30\textwidth}
% element: 03-s020:e03
Potreban je i izlaz \textcolor{DEisticanje}{\textbf{EO}} (\emph{enable output}),
koji dozvoljava koderu \textbf{nižeg prioriteta} da koduje.
\par\vspace{2mm}
Dodavanjem EI i EO osnovnoj šemi dobijamo prikazani koder prioriteta.
\end{column}
\end{columns}
\end{frame}

% Original: PDF strana 10, donji slajd; drugi deo.
\slajdid{03-s020b}
\begin{frame}[t,label=03-s020b]{Prenos dozvole signalom EO}
\setlength{\parskip}{0pt}
\fontsize{11}{13}\selectfont
% element: 03-s020b:e04
Na prvi pogled bismo umesto EO mogli koristiti invertovani GS.
Međutim, dozvola mora da se prenese \textbf{od višeg ka nižem prioritetu},
bez obzira na stanje pojedinačnog kodera.
\par\vspace{2mm}
\begin{columns}[T,onlytextwidth]
\begin{column}{.53\textwidth}
\textcolor{DEisticanje}{\textbf{Zašto invertovani GS nije dovoljan?}}
\par\vspace{1.5mm}
Uslovljeni GS može biti nula zato što je koderu \textbf{zabranjen rad}.
Njegova inverzija tada bi dozvolila nižem koderu da radi, \textbf{a ne sme}.
\par\vspace{1.5mm}
GS možemo ostaviti i neuslovljen: to ne smeta drugom nivou kodera prioriteta.
Ipak, ni on ne može zameniti EO. To što ova grupa nema šta da koduje
ne znači da \textbf{grupe višeg prioriteta nisu imale zahtev}.
\end{column}
\begin{column}{.44\textwidth}
\textcolor{DEisticanje}{\textbf{Suština izlaza EO}}
\par\vspace{1.5mm}
\textbf{Dozvoljeno mi je; nemam šta da kodujem:}
dozvoljavam kodovanje nižem koderu.
\par\vspace{2mm}
\textbf{Dozvoljeno mi je; imam šta da kodujem:}
zabranjujem kodovanje nižim koderima.
\par\vspace{2mm}
\textbf{Nije mi dozvoljeno da kodujem:}
ne dozvoljavam ni nižim koderima da koduju.
\end{column}
\end{columns}
\end{frame}

% Original: PDF strana 11, gornji slajd; prvi deo.
\slajdid{03-s021}
\begin{frame}[t,label=03-s021]{Lanac kodera prioriteta sa 16 ulaza}
\setlength{\parskip}{0pt}
% element: 03-s021:e02
Sa ovako realizovanim koderom prioriteta mreža sa \textbf{16 ulaza} postaje:
\par\vspace{1mm}
\begin{columns}[c,onlytextwidth]
\begin{column}{.47\textwidth}\centering
% element: 03-s021:e01
\begin{tikzpicture}[circuit logic US,x=.82mm,y=.67mm,font=\fontsize{7.8}{9}\selectfont,
 zbir/.style={or gate US,draw,line width=.8pt,logic gate inputs=nnnn,minimum width=4.8mm,minimum height=5mm,scale=.52}]
\foreach \g/\yy in {3/54,2/36,1/18,0/0}{
 \begin{scope}[shift={(0,\yy)}]
 \draw[DEvod] (0,0) rectangle (16,16);
 \node at (8,9) {KP};
 \foreach \n/\yin in {3/12.5,2/9.4,1/6.3,0/3.2}{
  \pgfmathtruncatemacro{\ii}{4*\g+\n}
  \node[anchor=west,inner sep=.5pt] at (0,\yin) {I\n};
  \draw[DEvod] (-2,\yin)--(0,\yin);
  \node[anchor=east,inner sep=.5pt] at (-2,\yin) {I\ii};
 }
 \foreach \port/\yout/\lab in {gs/10.8/GS,a1/6.5/A1,a0/3.2/A0}{
  \node[anchor=east,inner sep=.5pt] at (16,\yout) {\lab};
  \draw[DEvod] (16,\yout)--(18,\yout);
  \coordinate (p\g-\port) at (18,\yout);
 }
 \node[anchor=north,inner sep=.5pt] at (8,16) {EI};
 \draw[DEvod] (8,16)--(8,17);
 \coordinate (p\g-ei) at (8,17);
 \node[anchor=south,inner sep=.5pt] at (8,0) {EO};
 \draw[DEvod] (8,0)--(8,-1);
 \coordinate (p\g-eo) at (8,-1);
 \end{scope}
}
\draw[DEvod,DEisticanje] (p3-ei)--++(-3,0) node[left,inner sep=.5pt] {1};
\foreach \g/\next in {3/2,2/1,1/0}
 \draw[DEvod,DEisticanje] (p\g-eo)--(p\next-ei);
\draw[DEvod] (47,27) rectangle (63,43);
\node at (55,36) {KP};
\foreach \n/\yin in {3/39.5,2/36.4,1/33.3,0/30.2}
 \node[anchor=west,inner sep=.5pt] at (47,\yin) {I\n};
\foreach \g/\trunk/\yin in {3/43/39.5,2/41/36.4,1/37/33.3,0/39/30.2}
 \draw[DEvod] (p\g-gs)--(\trunk,0 |- p\g-gs)|-(47,\yin);
\node[anchor=north,inner sep=.5pt] at (55,43) {EI};
\draw[DEvod,DEisticanje] (55,43)--(55,45) node[above,inner sep=.5pt] {1};
\foreach \port/\yy/\inner/\outer in {gs/37.8/GS/GS,a1/33.5/A1/A3,a0/30.2/A0/A2}{
 \node[anchor=east,inner sep=.5pt] at (63,\yy) {\inner};
 \draw[DEvod] (63,\yy)--(65,\yy) node[right,inner sep=.5pt] {\outer};
}
\node[anchor=south,inner sep=.5pt] at (55,27) {EO};
\draw[DEvod] (55,27)--(55,25)--(65,25) node[right,inner sep=.5pt] {EO};
\node[zbir] (a1) at (40,-1) {};
\node[zbir] (a0) at (40,-10.5) {};
\foreach \g/\pin/\xx in {3/1/31,2/2/27,1/3/23,0/4/19}
 \draw[DEvod] (p\g-a1)--(\xx,0 |- p\g-a1)|-(a1.input \pin);
\foreach \g/\pin/\xx in {3/1/33,2/2/29,1/3/25,0/4/21}
 \draw[DEvod] (p\g-a0)--(\xx,0 |- p\g-a0)|-(a0.input \pin);
\foreach \name/\lab in {a1/A1,a0/A0}
 \draw[DEvod] (\name.output)--++(2,0) node[right,inner sep=.5pt] {\lab};
\end{tikzpicture}

\end{column}
\begin{column}{.50\textwidth}
% element: 03-s021:e03
Mrežu možemo dodatno pojednostaviti: izlaze \(A_1,A_0\) neaktivnog kodera
postavimo u \textcolor{DEisticanje}{\textbf{stanje visoke impedanse}}.
\par\vspace{2mm}
Tada odgovarajuće izlaze možemo \textbf{direktno spojiti},
pa nam \textcolor{DEisticanje}{\textbf{četvoroulazni ILI gejtovi nisu potrebni}}.
\end{column}
\end{columns}
\end{frame}

% Original: PDF strana 11, gornji slajd; drugi deo.
\slajdid{03-s021b}
\begin{frame}[t,label=03-s021b]{Crtanje simbola kola srednjeg stepena integracije}
\setlength{\parskip}{0pt}
% element: 03-s021b:e04
Verovatno ste uočili da nema strogog standarda za crtanje simbola
kola srednjeg stepena integracije.
\par\vspace{3mm}
\begin{columns}[c,onlytextwidth]
\begin{column}{.36\textwidth}\centering
\Oznaka{Pravougaoni simbol komponente}
\par\vspace{1mm}
\begin{tikzpicture}[x=1mm,y=1mm,font=\fontsize{9}{11}\selectfont]
\draw[DEvod] (0,0) rectangle (28,32);
\node at (14,23) {KP};\node at (14,18) {4/2};
\foreach \lab/\yy in {I3/26,I2/21,I1/16,I0/11,EI/6}{
 \node[anchor=west,inner sep=1pt] at (0,\yy) {\lab};
 \draw[DEvod] (-3,\yy)--(0,\yy);
}
\foreach \lab/\yy in {GS/26,A1/20,A0/14,EO/6}{
 \node[anchor=east,inner sep=1pt] at (28,\yy) {\lab};
 \draw[DEvod] (28,\yy)--(31,\yy);
}
\end{tikzpicture}

\end{column}
\begin{column}{.60\textwidth}
\textcolor{DEisticanje}{\textbf{Uobičajeno je da svi ulazi budu na jednoj strani
simbola, a svi izlazi na drugoj.}}
\par\vspace{3mm}
\textbf{Ulaze i izlaze ne mešamo na istoj stranici simbola.}
\par\vspace{3mm}
Dosta često crtamo \textbf{ulaze sa leve}, a \textbf{izlaze sa desne strane}.
\end{column}
\end{columns}
\end{frame}

\end{document}
